\input{../tex-inputs/latex-leseansicht-vorspann}

\section[Gustav Schwarzkopf an Arthur Schnitzler, 12. 7. 1900]{L04304 Gustav Schwarzkopf an Arthur Schnitzler, 12. 7. 1900}
\nopagebreak\mylabel{L04304v}
\rehead{ }\normalsize\beginnumbering\briefempfaengerindex{Schnitzler, Arthur@\textsc{Schnitzler, Arthur}!zzzSchwarzkopf, Gustav@\emph{von Gustav Schwarzkopf}!1900-07-122@{12. 7. 1900}|(be}
\toendnotes[C]{\smallbreak\pagebreak[2]}
\correspDesc{Versand  durch Gustav Schwarzkopf am 12. 7. 1900 in Wien
\newline{}Erhalt  durch Arthur Schnitzler im Zeitraum [13. 7. 1900
                  – 17. 7. 1900?] in Reichenau an der Rax}\toendnotes[C]{\smallbreak}
\Standort{CUL, Schnitzler, B 96a.}
\physDesc{Briefkarte, 1089 Zeichen
\newline{}Handschrift: schwarze Tinte, deutsche Kurrent}\toendnotes[C]{\smallbreak}
\pstart
           \raggedleft{}{\pb}12. 7. 00\pend
           \vspace{0.5em}
\pstart
           Lieber Arthur,{ }\label{K_L04304-1v}\edtext{Karte}{\lemma{\textnormal{\emph{Karte}}}\Cendnote{\textnormal{XXXX Auszeichnungsfehler: Dokument L04083 nicht gefunden. }}}\label{K_L04304-1}, Novelle\pwindex{Schnitzler, Arthur 15.\,5.\,1862 Wien – 21.\,10.\,1931 ebd.@\textsc{Schnitzler, Arthur} (15.\,5.\,1862 Wien – 21.\,10.\,1931 ebd.), \emph{Schriftsteller, Mediziner}!Frau Bertha Garlan. Roman@\strich\emph{Frau Bertha Garlan. Roman}|pwv} und \label{K_L04304-2v}\edtext{Brief}{\lemma{\textnormal{\emph{Brief}}}\Cendnote{\textnormal{XXXX Auszeichnungsfehler: Dokument L04080 nicht gefunden. }}}\label{K_L04304-2} erhalten, Ihrem Wunſch
               gemäß verſpare ich die Äußerungen über die Novelle\pwindex{Schnitzler, Arthur 15.\,5.\,1862 Wien – 21.\,10.\,1931 ebd.@\textsc{Schnitzler, Arthur} (15.\,5.\,1862 Wien – 21.\,10.\,1931 ebd.), \emph{Schriftsteller, Mediziner}!Frau Bertha Garlan. Roman@\strich\emph{Frau Bertha Garlan. Roman}|pwv} auf mündliche Unterhaltung, obwol ich faſt nur
               Gutes zu{ }ſagen hätte, was auch{ }ſchriftlich leichter zu erledigen iſt. Aber man ka{\geminationn} ja wirklich in 5 Minuten mehr darüber{ }ſprechen, als in
               vier Bogen{ }ſchreiben. Nach Reichenau\oindex{Reichenau an der Rax@\textbf{Reichenau an der Rax}, \emph{Verwaltungsgebiet}|pw} wäre ich{ }ſchon geko{\geminationm}en, we{\geminationn} Sie
               nicht gerade \uline{jetzt} dort wären. Jetzt muß ich in Wien\oindex{Wien@\textbf{Wien}, \emph{Verwaltungsgebiet}|pw} bleiben, weil Max\pwindex{Schwarzkopf, Max 12.\,6.\,1857 Wien – 14.\,4.\,1928 ebd.@\textsc{Schwarzkopf, Max} (12.\,6.\,1857 Wien – 14.\,4.\,1928 ebd.), \emph{Rechtsanwalt}|pw} vorgeſtern abgereiſt iſt; habe auch – unter dieſer
               Motivierung – eine{ }ſehr liebenswürdige Einladung zu Ganghofer\pwindex{Ganghofer, Ludwig 7.\,7.\,1855 Kaufbeuren – 24.\,7.\,1920 Tegernsee@\textsc{Ganghofer, Ludwig} (7.\,7.\,1855 Kaufbeuren – 24.\,7.\,1920 Tegernsee), \emph{Schriftsteller}|pw}\pwindex{Ganghofer, Katharina 7.\,7.\,1856 Pest – 8.\,4.\,1930 München@\textsc{Ganghofer, Katharina} (7.\,7.\,1856 Pest – 8.\,4.\,1930 München), \emph{Schauspielerin}|pw}{ }{\pb}abgelehnt, die ich vor einer Woche
               erhalten. Auf einen Tag kö{\geminationn}te ich wol leicht fort, aber
                  da{\geminationn} iſt da{\geminationn} doch die
               Fahrt zu thuen und da Sie in einer Woche doch noch nach Wien\oindex{Wien@\textbf{Wien}, \emph{Verwaltungsgebiet}|pw} ko{\geminationm}en. – Hier iſt’s im
                  \textcolor{gray}{Durch}haus noch recht belebt, das Wetter trägt dazu bei, daß man
               glauben ka{\geminationn} im Herbſt zu{ }ſein. Eberma{\geminationn}\pwindex{Ebermann, Leo 16.\,7.\,1863 Draganovka – 9.\,10.\,1914 Wien@\textsc{Ebermann, Leo} (16.\,7.\,1863 Draganovka – 9.\,10.\,1914 Wien), \emph{Schriftsteller, Journalist, Rechtswissenschaftler}|pw}, Kapper\pwindex{Kapper, Hans 10.\,3.\,1891 Wien – September 1962 New York City@\textsc{Kapper, Hans} (10.\,3.\,1891 Wien – September 1962 New York City), \emph{Schriftsteller}|pw}, der kleine Kraus\pwindex{Kraus, Karl 28.\,4.\,1874 Jičín – 12.\,6.\,1936 Wien@\textsc{Kraus, Karl} (28.\,4.\,1874 Jičín – 12.\,6.\,1936 Wien), \emph{Schriftsteller, Publizist, Schriftsteller}|pw}, auch \textsc{Ludaſsy\pwindex{Gans-Ludassy, Julius von 13.\,4.\,1858 Wien – 30.\,9.\,1922 ebd.@\textsc{Gans-Ludassy, Julius von} (13.\,4.\,1858 Wien – 30.\,9.\,1922 ebd.), \emph{Schriftsteller, Journalist, Herausgeber}|pw}}, der Reklamefeind, zeigt sich. (Der Arme. Die Notizen über »d. letzte Knopf\pwindex{Gans-Ludassy, Julius von 13.\,4.\,1858 Wien – 30.\,9.\,1922 ebd.@\textsc{Gans-Ludassy, Julius von} (13.\,4.\,1858 Wien – 30.\,9.\,1922 ebd.), \emph{Schriftsteller, Journalist, Herausgeber}!letzte Knopf. Volksstück in drei Aufzügen@\strich\emph{Der letzte Knopf. Volksstück in drei Aufzügen}|pw}« verfolgen ihn unbarmherzig auch noch im
                  Juli!– –\pend
           
\pstart
           Bitte, empfehlen Sie mich beſtens Ihrer Frau Mama\pwindex{Schnitzler, Louise 8.\,7.\,1840 Kőszeg – 9.\,9.\,1911 Wien@\textsc{Schnitzler, Louise} (8.\,7.\,1840 Kőszeg – 9.\,9.\,1911 Wien)|pwv}, Schweſter\pwindex{Hajek, Gisela 20.\,12.\,1867 Wien – 3.\,2.\,1953 Cambridge@\textsc{Hajek, Gisela} (20.\,12.\,1867 Wien – 3.\,2.\,1953 Cambridge)|pwv} und Schwägerin\pwindex{Schnitzler, Helene 16.\,7.\,1871 Budapest – September 1941 Atlantischer Ozean@\textsc{Schnitzler, Helene} (16.\,7.\,1871 Budapest – September 1941 Atlantischer Ozean)|pwv},\pend
           
\pstart
           Herzlichſt{\\[\baselineskip]}Ihr{\\[\baselineskip]}\spacefill\mbox{Gustav.}\pend
           \leftskip=0em{}\selectlanguage{ngerman}\endnumbering\briefempfaengerindex{Schnitzler, Arthur@\textsc{Schnitzler, Arthur}!zzzSchwarzkopf, Gustav@\emph{von Gustav Schwarzkopf}!1900-07-122@{12. 7. 1900}|)be}\mylabel{L04304h}
\begin{anhang}
\end{anhang}\newcommand{\dateiname}{L04304}\newcommand{\titel}{Gustav Schwarzkopf an Arthur Schnitzler, 12. 7. 1900}\newcommand{\editorInnen}{Herausgegeben von Jahnke, SelmaMüller, Martin Anton}\input{../tex-inputs/latex-leseansicht-abspann}
